\section{Datalake Structure Evaluation}
\label{sec:datalake}

The datalake stores the same header/body pairs under three layouts. This
section evaluates them on the dimensions the assignment asks for: write and
download throughput, document lookup, discovery of new books, storage
overhead, and recovery after an interruption. All numbers come from
\texttt{datalake\_experiments.py} (Python baseline; the network probe from
\texttt{download\_probe.py}) and the equivalent
\texttt{datalake\_bench} code in each port, with the parameters of
\cref{tab:workload}.

\subsection{Write and download throughput}
\label{sec:datalake-throughput}

Write throughput is one timed run per layout that reads each of the
\fixed{100} tracked headers and bodies from \texttt{data\_source/} and stores
them through the layout's storage path. Since the current snapshot every
implementation writes each file to a temporary sibling, forces it to stable
storage with \texttt{fsync} and renames it into place
(\cref{sec:refactor-effect}). The download is a separate experiment:
\fixed{10} books fetched sequentially in each of \fixed{5} rounds through one
pooled connection per language, timed from before the HTTP request until the
body is written to disk. The rounds interleave the languages round-robin, so
no language downloads at a systematically different hour of the run. Both
experiments operate on the same \fixed{100}-book corpus.

\paragraph{Writes.}
Go stored the corpus at 46.9--48.1 items/s, Rust at 38.3--39.0, Java at
30.5--35.1 and Python at 26.4--29.0 (\cref{fig:throughput}a): 21~ms per book
for Go and 36~ms for Python, each book being two synchronised files. The
layout matters little. Within each implementation the three layouts differ
by at most 3\% in Go and Rust, 10\% in Python and 15\% in Java, and the
book-based layout, which creates one directory per book (plus, in Python, a
\texttt{metadata.json} per book), is the slowest in Python, Go and Rust.
Client CPU utilisation during the writes is 10--13\% (Rust), 17--18\% (Go)
and 27--37\% (Java): the phase is bound by storage synchronisation, not by
computation. The remaining ordering of the ports (Go, Rust, Java) is not
explained by the data collected.

\paragraph{Downloads.}
\Cref{tab:download} summarises the five rounds. The probe is network-bound
(client CPU 1.5--6.3\% in Go and Rust), so it characterises the HTTP clients
and the path to Project Gutenberg rather than computation. Python
(3.24~items/s) and Go (3.12) overlap within one standard deviation; Rust is
lower (2.60) and Java lowest (1.88). Java's mean is pulled down by a first
round of 11.2~s (0.89~items/s; median 2.13), and its client uses 25--44\% CPU;
since every round starts a fresh JVM through Maven, start-up and JIT cost is
a plausible cause, which we did not isolate. Panel~(b) of
\cref{fig:throughput} shows the fifth round only, not these means
(\cref{sec:radar}).

\begin{table}[htbp]
    \centering
    \caption{Download + write throughput, \fixed{10} books per round,
    \fixed{5} interleaved rounds per language (\texttt{throughput.csv}).
    The previous snapshot had a single round per language, taken at
    different hours and without connection pooling in Python and Rust.}
    \label{tab:download}
    \begin{tabular}{@{}lcccc@{}}
        \toprule
        Language & Mean $\pm$ SD (items/s) & Min--max & Median & Previous ($n=1$) \\
        \midrule
        Python & 3.24 $\pm$ 0.29 & 2.90--3.52 & 3.39 & 1.48 \\
        Go & 3.12 $\pm$ 0.37 & 2.67--3.47 & 3.31 & 3.32 \\
        Rust & 2.60 $\pm$ 0.25 & 2.25--2.96 & 2.64 & 1.32 \\
        Java & 1.88 $\pm$ 0.56 & 0.89--2.22 & 2.13 & 2.87 \\
        \bottomrule
    \end{tabular}
\end{table}

\subsection{Document lookup}
\label{sec:datalake-lookup}

\fixed{100} books are located once per layout: the harness resolves the
book's directory and checks that both its header and its body exist; the
files are not read. The book-based and batch-based layouts derive the
directory arithmetically from the ID, whereas the time-based layout must list
its hour directories, newest first, and probe each until it finds the book.
\Cref{tab:datalake-lookup} reports the mean and, because several series are
right-skewed, the minimum.

\begin{table}[htbp]
    \centering
    \caption{Datalake lookup latency, 100 books per layout, in $\mu$s: mean
    (minimum) per language; raw rows in
    \texttt{benchmarks\_results/<language>/statistics.csv}.}
    \label{tab:datalake-lookup}
    \begin{tabular}{@{}lcccc@{}}
        \toprule
        Layout & Python & Go & Rust & Java \\
        \midrule
        Time-based & 90 (76) & 25 (20) & 19 (17) & 433 (103) \\
        Book-based & 39 (32) & 7 (6) & 7 (5) & 14 (8) \\
        Batch-based & 35 (30) & 7 (6) & 6 (5) & 50 (9) \\
        \bottomrule
    \end{tabular}
\end{table}

The time-based layout is the slowest in every implementation (19--433~$\mu$s
mean, 17--103~$\mu$s minimum), as its directory listing predicts. Book- and
batch-based lookups are indistinguishable in Python, Go and Rust (means
within 4~$\mu$s, minima within 2~$\mu$s). Java's means are dominated by a few
slow samples (standard deviations of 44--1{,}294~$\mu$s, maxima up to
11.7~ms): this phase runs without warm-up (\cref{sec:warmup}), so its first
lookups include interpretation and JIT compilation; on the minima Java
follows the same pattern as the other three. Every mean stays below
0.5~ms.

\subsection{Discovery of downloaded-but-not-indexed books}
\label{sec:datalake-discovery}

Lookup answers ``where is a book I know exists''. The complementary
operation is discovery: \emph{which books are stored but not yet indexed?}
The experiment marks the first \fixed{50} of the \fixed{100} stored IDs as
already indexed, in an in-memory set, and times one scan per layout that
lists the stored IDs and returns the 50 that are not in the set
(\texttt{incremental\_<layout>} in \texttt{benchmarks.csv}). The initial
state is identical across layouts and languages; the control files of
\cref{sec:architecture} are not involved and no book is added out of band.
\Cref{tab:discovery} lists the single sample per layout and language.

\begin{table}[htbp]
    \centering
    \caption{Discovery scan, 100 stored books of which 50 are marked
    indexed, in ms (one sample per cell).}
    \label{tab:discovery}
    \begin{tabular}{@{}lcccc@{}}
        \toprule
        Layout & Python & Go & Rust & Java \\
        \midrule
        Time-based & 2.80 & 0.76 & 2.04 & 22.55 \\
        Book-based & 5.39 & 1.98 & 3.23 & 18.48 \\
        Batch-based & 2.94 & 0.76 & 0.91 & 3.38 \\
        \bottomrule
    \end{tabular}
\end{table}

The scan took 0.76--22.6~ms. The batch-based layout is the fastest or tied
in all four implementations, and the book-based layout, whose scan lists 100
directories instead of 2--3, is the slowest in Python, Go and Rust. With one
sample per cell these are observations, not a ranking. The CPU figures
recorded for these few-millisecond rows reach 489\% despite single-core
pinning, which shows that process CPU time is accounted at a coarser
resolution than the operation lasts; \texttt{cpu\_percent} is therefore not
interpreted for operations shorter than about 10~ms.

\subsection{Storage overhead}
\label{sec:datalake-storage}

After storing the identical corpus, each layout is measured for apparent
bytes (the sum of file sizes), file count and directory count
(\texttt{disk\_usage.csv}; the harness's ``MB'' are MiB, $2^{20}$ bytes).
The layouts differ only in how many intermediate directories they create: a
time-based layout adds a day and an hour level, a book-based layout one
directory per book, a batch layout one directory per range of 1{,}000 IDs.
\Cref{tab:datalake-storage} reports the figures for the shared
\fixed{100}-book corpus.

\begin{table}[htbp]
    \centering
    \caption{On-disk cost of the same corpus under each layout (apparent
    size of the header and body files; allocated blocks were not measured).
    Ranges span the four implementations: Python stores 1.29~MiB less
    because it reads the corpus with universal newlines and drops carriage
    returns, and adds a \texttt{metadata.json} per book in the book-based
    layout.}
    \label{tab:datalake-storage}
    \begin{tabular}{@{}lccc@{}}
        \toprule
        Layout & Size (MiB) & Files & Directories \\
        \midrule
        Time-based & 57.38--58.67 & 200 & 2 \\
        Book-based & 57.40--58.67 & 200--300 & 100 \\
        Batch-based & 57.38--58.67 & 200 & 3 \\
        \bottomrule
    \end{tabular}
\end{table}

The directory counts are consistent with the two measured differences
between layouts: the book-based layout's 100 directories go with the slowest
discovery scan (\cref{tab:discovery}), and the time-based layout's directory
listing with the slowest lookup (\cref{tab:datalake-lookup}). The counts grow
differently with the corpus (one directory per book, per 1{,}000 IDs, or per
ingestion hour), but only the 100-book state was measured.

\subsection{Recovery after interruption}
\label{sec:datalake-recovery}

What the harness measures is a \emph{resume-detection check}, not a crash
test. With all \fixed{100} books already stored by the write phase, the first
\fixed{50} IDs are marked as processed in memory; the harness then times
(i)~the detection of the 50 pending books, with the same scan as discovery,
and (ii)~locating and reading back their headers and bodies
(\texttt{recovery.csv}). No process is killed, the control files are not
used and no partial write can occur, so the reported 0 duplicated and 0 lost
books follow from the design and are \emph{not} evidence of crash safety.
The durable-write path now shared by all four implementations (temporary
file, \texttt{fsync}, rename) means a crashed writer leaves at most a
leftover temporary file, never a truncated book, but this experiment does
not exercise that property; a kill-and-restart test is listed as future work
(\cref{sec:future}).

\begin{table}[htbp]
    \centering
    \caption{Resume-detection check: 100 books stored, the first 50 marked
    processed in memory, one sample per layout and language. Duplicated and
    lost counts are zero by construction. Times are min--max across the four
    implementations, in ms.}
    \label{tab:recovery}
    \begin{tabular}{@{}lcccccc@{}}
        \toprule
        Layout & Before & After & Duplicated & Lost & Detection (ms) & Total (ms) \\
        \midrule
        Time-based & 50 & 50 & 0 & 0 & 0.70--10.79 & 15.9--96.7 \\
        Book-based & 50 & 50 & 0 & 0 & 1.61--14.72 & 18.1--102.5 \\
        Batch-based & 50 & 50 & 0 & 0 & 0.81--4.07 & 15.7--80.2 \\
        \bottomrule
    \end{tabular}
\end{table}
